\documentclass[12pt]{article}
\usepackage[a4paper, margin=2.5cm]{geometry}
\usepackage{amsmath}
\usepackage{booktabs}
\usepackage{caption}
\usepackage{hyperref}
\usepackage{setspace}
\usepackage{microtype}
\usepackage{lmodern}
\usepackage[T1]{fontenc}
\usepackage[utf8]{inputenc}
\usepackage{siunitx}
\sisetup{group-separator={,}, group-minimum-digits=4}

\onehalfspacing

\title{Data Acquisition for Equity Index Options:\\
S\&P~500, Euro~Stoxx~50, S\&P/ASX~200, and Swiss Market Index}
\author{}
\date{April 2026}

\begin{document}

\maketitle

\noindent
This note describes in precise detail the construction of a historical panel of
equity index option contracts and their associated price time series for four major
index option markets: the S\&P~500 (SPX, CBOE), the Euro~Stoxx~50 (STOXX50E,
Eurex), the S\&P/ASX~200 (AXJO, ASX), and the Swiss Market Index (SMI, Eurex).
All data are sourced exclusively through the Refinitiv Eikon API,
specifically the \texttt{ek.get\_data()} endpoint of the \texttt{eikon} Python
package. No alternative retrieval methods---such as the \texttt{rd} library or
\texttt{ek.get\_timeseries()}---are employed, as these offer a narrower set of
historical fields and less reliable coverage of expired contracts. The Eikon
Desktop application must be active on the local machine at the time of any data
fetch, as the API communicates via a local proxy process (default port 9000).

The overarching objective is to compile, for each index, the complete strip of
monthly European-style options traded on or before a given reference date,
together with their full bid--ask price history up to expiration. The data are
stored in a PostgreSQL relational database under the schema \texttt{pub\_options},
using two primary tables: \texttt{options\_chains}, which records contract-level
metadata, and \texttt{options\_prices}, which records the daily bid and ask prices
of each contract. The primary key of \texttt{options\_chains} is the pair
(\texttt{underlying\_ric},~\texttt{option\_ric}), where \texttt{option\_ric}
stores the Eikon base-form identifier of the contract, free of any leading slash
or historical suffix. The primary key of \texttt{options\_prices} is
(\texttt{option\_ric},~\texttt{price\_date}), where the \texttt{option\_ric}
field in this table carries the historical identifier including the
\texttt{\^{}} suffix described below.

\bigskip

\noindent\textbf{Instrument Identification and RIC Structure.}
Each option contract in the Eikon system is uniquely identified by a Reuters
Instrument Code (RIC). The structure of these codes differs substantially across
exchanges, and a precise understanding of the encoding is essential for
systematic retrieval. For the S\&P~500, the CBOE employs a
date-embedded, strike-scaled format. A SPX option is identified as
\texttt{SPX\{ltr\}\{DD\}\{YY\}\{K\}.U}, where \texttt{ltr} is a lowercase
single character encoding both the option type and calendar month (a--l for
January through December calls; m--x for January through December puts),
\texttt{DD} is the two-digit expiration day, \texttt{YY} is the two-digit
calendar year, \texttt{K} is the strike price multiplied by 10 and
zero-padded to five digits, and \texttt{.U} denotes the CBOE venue. As a
concrete example, the SPX February~2026 call at a strike of 6{,}900 is
identified as \texttt{SPXb202669000.U}, in which \texttt{b} signals a February
call, \texttt{20} is the expiration day, \texttt{26} encodes 2026, and
\texttt{69000} encodes a strike of \$6{,}900.

For the Euro~Stoxx~50, Eurex uses the format
\texttt{STXE\{K{\times}10\}\{ltr\}\{y\}.EX}, where the strike is encoded as ten
times the nominal strike value (so that a strike of 5{,}000 appears as
\texttt{50000}), \texttt{ltr} is a single uppercase month letter (A for January
through L for December for calls, M through X for puts), \texttt{y} is a
single-digit calendar year, and \texttt{.EX} denotes the Eurex venue. A
January~2026 Euro~Stoxx~50 call at strike 5{,}000 is thus encoded as
\texttt{STXE50000A6.EX}. The SMI options traded on Eurex follow an analogous
convention under the prefix \texttt{OSMI}, with the same strike multiplication
by ten and single-digit year, using \texttt{.EX} as the venue suffix. A
January~2026 SMI call at strike 12{,}000 appears as \texttt{OSMI120000A6.EX}.
For the S\&P/ASX~200, the ASX exchange uses the format
\texttt{AXJO\{K\}\{ltr\}\{y\}.AX}, where the strike is recorded at face value
without any multiplication, \texttt{ltr} and \texttt{y} follow the same
uppercase-letter and single-digit-year convention, and \texttt{.AX} is the
venue suffix. An ASX~200 April~2026 put at strike 7{,}500 reads as
\texttt{AXJO7500P6.AX}, with \texttt{P} denoting the April put letter.

A structurally important feature of the single-digit year encoding used by
Eurex and ASX is that the RIC for a given contract repeats every ten years.
The identifier \texttt{STXE50000A6.EX} is syntactically identical for
January~2016 and January~2026 contracts. Eikon resolves the ambiguity by
serving data for the currently active cycle: as of April~2026, queries for
\texttt{STXE50000A6.EX} return January~2026 data, not January~2016. For
all historical price retrieval of expired contracts, Eikon instead requires a
disambiguating suffix of the form \texttt{\^{}\{LETTER\}\{YY\}}, appended
directly to the base RIC. The letter in the suffix is always the call-series
month letter, regardless of whether the underlying contract is a call or a put,
and \texttt{YY} is the full two-digit calendar year. Thus the historical
identifier for a Euro~Stoxx~50 December~2025 put at strike 4{,}000 is
\texttt{STXE40000X5.EX\^{}L25}, where the suffix \texttt{\^{}L25} uses
\texttt{L} (the December call letter) and the two-digit year 25. The SPX,
despite already encoding a two-digit year within its base RIC format, follows
the same suffix convention for historical retrieval of expired contracts.
The column \texttt{option\_ric\_hist} in \texttt{options\_chains} stores this
historical identifier for every contract, and it is this column---not
\texttt{option\_ric}---that is used as the join key to \texttt{options\_prices}.
A live query against an active contract uses the format
\texttt{/\{option\_ric\}} (with a leading slash), while historical time-series
retrieval uses the \texttt{option\_ric\_hist} form without a leading slash.

\bigskip

\noindent\textbf{Monthly Expiry Convention and Expiry Classification.}
European index options on the four markets expire on exchange-specific calendar
dates. For SPX and the two Eurex products (Euro~Stoxx~50 and SMI), the standard
monthly expiration is the third Friday of the calendar month, adjusted backward
by one business day when that Friday coincides with a public holiday. Notable
exceptions include Good Friday falling on the third Friday of April in 2019,
2022, and 2025, in each of which the expiration shifted to the preceding
Thursday. The ASX convention for the S\&P/ASX~200 differs: monthly expirations
fall on the third Thursday of each month. Each contract in \texttt{options\_chains}
is assigned an \texttt{expiry\_type} field taking one of three values: \texttt{M}
(standard monthly), \texttt{W} (weekly, i.e.\ any expiration other than the
canonical monthly date), or \texttt{L} (long-dated, defined as options with
time to expiration exceeding one year at the time of the monthly pipeline
snapshot, applicable to equity options only). The present dataset is
restricted almost entirely to monthly expirations; weekly series account for
less than one percent of contracts across all four indices and arise only
because exchange chain queries occasionally return near-term weekly contracts
listed alongside the regular monthly series.

\bigskip

\noindent\textbf{Strike Candidate Generation and Chain Construction.}
The construction of the option strip proceeds through a systematic candidate
generation step applied independently to each (index, expiry month) pair.
For each target month, the index closing level on or immediately before the
expiry date is retrieved from the \texttt{pub\_equity.index\_prices\_daily} table
(with a fallback Eikon query covering the preceding seven calendar days).
Candidate strikes are then generated as the arithmetic sequence of all
multiples of five index points within a symmetric band of $\pm60\%$ around the
observed index level, subject to a floor of one strike step. Concretely, for
an index level $S$, the candidate set is
\[
  \mathcal{K} = \bigl\{\, k \in 5\mathbb{Z}^+ \colon
  \lfloor 0.4S/5 \rfloor \cdot 5 \;\le\; k \;\le\;
  \lceil 1.6S/5 \rceil \cdot 5 \,\bigr\},
\]
so that at the typical index levels observed over the sample period, each
monthly expiry for SPX generates approximately 335 distinct strike levels on
each of the call and put sides. The equivalent figures for Euro~Stoxx~50, SMI,
and S\&P/ASX~200 are approximately 130, 82, and 110 strike levels per side,
respectively, reflecting differences in index levels and the implied width of
the $\pm60\%$ band in points.

For each candidate (option type, strike, expiry) triple, an Eikon RIC is
constructed according to the exchange-specific encoding described above.
Contracts already present in the database are identified via a start-of-run
query on \texttt{options\_chains} and excluded from the fetch. Remaining
candidates are submitted to the Eikon \texttt{get\_data()} API in batches of 75
RICs per call, querying the fields \texttt{TR.BIDPRICE.Date},
\texttt{TR.BIDPRICE}, and \texttt{TR.ASKPRICE} over a lookback window of 18
months ending on the expiry date. A contract is considered to exist and is
written to \texttt{options\_chains} if and only if at least one non-null bid or
ask observation is returned. Non-existing strike levels---which are numerous
given the granularity of the five-point grid---simply return empty responses
and are discarded. The batch sleep between successive API calls is set to 0.15
seconds; upon encountering an HTTP~429 rate-limit error, the system suspends
execution for 3{,}600 seconds and retries, with up to 24 retry attempts per
batch, providing a theoretical maximum tolerance of 24 hours of sustained rate
limiting before a batch is abandoned to a secondary retry queue. The retry
queue is processed once the main pass is complete; only if a batch fails in
the retry queue as well is it flagged as a permanent failure, at which point
the recommended resume month is logged explicitly so that execution can be
restarted at the precise point of interruption.

\bigskip

\noindent\textbf{Price Time Series Acquisition.}
Once the chain metadata is populated in \texttt{options\_chains}, daily bid and
ask prices are retrieved separately via the monthly pipeline module
\texttt{fetch\_option\_prices}. For active (unexpired) contracts, the pipeline
queries Eikon using the live RIC format (\texttt{/\{option\_ric\}}) and
retrieves all available observations from the later of one year before the
pipeline snapshot date or the maximum previously recorded price date for that
contract. For expired contracts, the pipeline uses the historical
\texttt{option\_ric\_hist} identifier, targeting a window beginning 12 months
before the expiry date (or from the last recorded observation, if later).
Expired and active contracts are submitted simultaneously in mixed batches to
economise on API calls. Price records are written to \texttt{options\_prices}
with the \texttt{option\_ric\_hist} as the primary key component, so that
prices for both active and expired contracts share a unified, historically
unambiguous identifier. The join between the two tables is always
\texttt{options\_chains.option\_ric\_hist = options\_prices.option\_ric}.
A completeness flag (\texttt{bid\_complete}, \texttt{ask\_complete}) is
set to \texttt{TRUE} in \texttt{options\_chains} if at least one price
observation exists within 30 calendar days before the contract's expiry date,
providing a rapid diagnostic for data coverage.

\bigskip

\noindent\textbf{Descriptive Statistics.}
Table~\ref{tab:desc} summarises the composition of the option panel as of
1 April 2026. The sample spans a common window of January 2016 through March
2026 for monthly expirations, though additional long-dated SPX and Euro~Stoxx~50
expirations extend the chain metadata to December 2031 and December 2035,
respectively, reflecting contracts already listed on the exchanges at the time
of the most recent snapshot.

The SPX constitutes the largest subpanel with 93{,}513 individual option
records across 142 distinct expiry dates. The strike grid covers the range
from \$200 to \$16{,}000, encompassing 1{,}336 unique strike levels in total
across the full sample period, with an average of 334.7 strikes per monthly
expiry (median 345). The large variation in the per-expiry count---from a
minimum of 32 strikes to a maximum of 520---reflects the pronounced
time-variation in the index level over the decade-long sample, which spans the
COVID-19 drawdown of early 2020, the subsequent recovery, and the sustained
rally through 2024--2025. The price panel for SPX comprises 13{,}262{,}266
daily bid--ask records drawn from 92{,}583 distinct contracts, equivalent to
a coverage rate of 99.0\% of all catalogued SPX options. The average bid price
across all non-zero daily observations is \$502.77 (median \$191.60), and the
average bid--ask spread is \$6.62.

The Euro~Stoxx~50 panel contains 38{,}659 option records across 150 expiry
dates, with strikes ranging from 100 to 18{,}000 index points. The average
number of distinct strike levels per monthly expiry is 130.1 (median 120),
reflecting the somewhat lower nominal index level relative to the $\pm60\%$
band width in absolute terms. The price panel encompasses 4{,}639{,}489
daily observations from 35{,}263 distinct contracts (coverage rate of
91.2\%). The average bid is \EUR{}435.18 with a median of \EUR{}180.10 and an
average bid--ask spread of \EUR{}13.34. The wider proportional spread relative
to SPX is consistent with the lower liquidity of Eurex-traded Euro~Stoxx~50
options at extreme strikes.

For the S\&P/ASX~200, the panel spans 131 monthly expiry dates and records
25{,}730 option contracts over strikes from A\$2{,}500 to A\$16{,}500.
Unlike the other three indices, no weekly or long-dated contracts appear in
the dataset, consistent with the ASX product structure that offers only
standard monthly expirations for index options. The average strike count
per monthly expiry is 109.8 (median 107). Price coverage is high at 94.7\%,
with 2{,}480{,}058 daily price records from 24{,}371 distinct contracts.
The average bid price is A\$395.18 (median A\$226.00) and the average
bid--ask spread is A\$38.85, substantially wider in nominal terms than the
SPX and Euro~Stoxx~50, which partly reflects the index's higher absolute
strike level relative to option premia and the market's lower depth.

The SMI panel contains 20{,}638 option records across 139 expiry dates,
with strikes ranging from CHF~3{,}500 to CHF~20{,}000 and an average of
82.0 strike levels per monthly expiry (median 81), the most concentrated
distribution of the four indices. The price panel covers 2{,}378{,}843
daily observations from 19{,}893 contracts (coverage rate of 96.4\%).
Average bid and ask prices are CHF~735.0 and CHF~774.0, respectively,
reflecting the high nominal index level; the median bid is CHF~377.3.
The average bid--ask spread of CHF~39.3 is comparable in magnitude to
the ASX~200.

The aggregate price panel across all four indices totals 22{,}760{,}656
daily bid--ask observations, providing a longitudinal record of option
pricing that is dense enough to support high-frequency cross-sectional
analyses of the volatility surface as well as time-series estimation
of model-free variance measures such as the synthetic volatility index
(SVIX) of \citet{martin2017} and its equity-level generalisation.

\begin{table}[ht]
\centering
\caption{Descriptive statistics of the option panel as of 1 April 2026.
``Contracts'' denotes the number of distinct (underlying, option RIC) pairs
in \texttt{options\_chains}. ``Expiries'' counts distinct expiry dates.
``Strikes per expiry'' refers to monthly expirations only.
``Price rows'' is the total number of daily bid--ask records in
\texttt{options\_prices}. Prices are denominated in the respective index
currency (USD for SPX; EUR for Euro~Stoxx~50; AUD for S\&P/ASX~200;
CHF for SMI). Spreads and price levels are computed over non-zero bid
observations.}
\label{tab:desc}
\small
\begin{tabular}{lrrrr}
\toprule
 & \textbf{SPX} & \textbf{Euro Stoxx 50} & \textbf{S\&P/ASX 200} & \textbf{SMI} \\
\midrule
Exchange & CBOE & Eurex & ASX & Eurex \\
Currency & USD & EUR & AUD & CHF \\
Expiry convention & 3rd Fri & 3rd Fri & 3rd Thu & 3rd Fri \\
Strike step (pts) & 5 & 5 & 5 & 5 \\
Strike range (\%) & $\pm60$ & $\pm60$ & $\pm60$ & $\pm60$ \\
\midrule
Contracts & 93{,}513 & 38{,}659 & 25{,}730 & 20{,}638 \\
\quad of which calls & 46{,}749 & 19{,}332 & 12{,}203 & 10{,}374 \\
\quad of which puts & 46{,}764 & 19{,}327 & 13{,}527 & 10{,}264 \\
\quad of which monthly & 93{,}093 & 38{,}393 & 25{,}730 & 20{,}452 \\
\quad of which weekly & 420 & 266 & 0 & 186 \\
Distinct expiry dates & 142 & 150 & 131 & 139 \\
First expiry & Jan 2016 & Jan 2016 & Jan 2016 & Jan 2016 \\
Last expiry (metadata) & Dec 2031 & Dec 2035 & Sep 2027 & Dec 2030 \\
Unique strike levels & 1{,}336 & 603 & 314 & 234 \\
Strike range & 200--16{,}000 & 100--18{,}000 & 2{,}500--16{,}500 & 3{,}500--20{,}000 \\
\midrule
Avg.\ strikes / monthly expiry & 334.7 & 130.1 & 109.8 & 82.0 \\
Median strikes / monthly expiry & 345 & 120 & 107 & 81 \\
Min.\ strikes / monthly expiry & 32 & 13 & 27 & 12 \\
Max.\ strikes / monthly expiry & 520 & 322 & 163 & 145 \\
\midrule
Price rows (total) & 13{,}262{,}266 & 4{,}639{,}489 & 2{,}480{,}058 & 2{,}378{,}843 \\
Price coverage (\%) & 99.0 & 91.2 & 94.7 & 96.4 \\
Price sample start & Jul 2014 & Dec 2014 & Feb 2015 & Dec 2014 \\
Price sample end & Apr 2026 & Apr 2026 & Apr 2026 & Apr 2026 \\
Rows with bid & 12{,}016{,}139 & 4{,}344{,}340 & 2{,}215{,}980 & 2{,}210{,}350 \\
Rows with ask & 12{,}016{,}993 & 4{,}355{,}703 & 2{,}237{,}517 & 2{,}213{,}674 \\
Avg.\ bid price & 502.77 & 435.18 & 395.18 & 735.0 \\
Median bid price & 191.60 & 180.10 & 226.00 & 377.3 \\
Avg.\ bid--ask spread & 6.62 & 13.34 & 38.85 & 39.3 \\
\bottomrule
\end{tabular}
\end{table}

\bibliographystyle{plainnat}
\begin{thebibliography}{1}
\bibitem{martin2017}
Martin, I. (2017).
\newblock What is the expected return on the market?
\newblock \textit{The Quarterly Journal of Economics}, 132(1):367--433.
\end{thebibliography}

\end{document}
